\documentclass[12pt]{article}
\usepackage[margin=1in]{geometry}
\usepackage{amsmath,amssymb}
\usepackage[hidelinks]{hyperref}
\usepackage{enumitem}
\setlist{nosep,leftmargin=*}

\title{In-Context Learning of Hidden Physical Dynamics from Push Interactions}
\author{David Jo \and Soomi \and \textit{[Other team members]}}
\date{CS 182 Final Project Proposal}

\begin{document}
\maketitle

\section*{Abstract}
Physical properties such as mass and friction cannot be determined from appearance alone; they must be inferred through interaction, and in-context learning (ICL) allows a frozen transformer to adapt from examples in its prompt.
Prior work suggests that transformers can infer dynamics in context, but it remains unclear which interactions make a context informative.
We train a small transformer from scratch to predict the displacement of a simulated cube with hidden mass and friction under a new force push, conditioned on previous pushes of the same cube.
We investigate how context length, diversity, redundancy, and measurement noise affect prediction accuracy, including on unseen mass and friction ranges.
We compare against an explicit Bayesian system-identification estimator and an oracle with access to the true parameters.
We aim to characterize when interaction history is sufficient for accurate in-context prediction and whether the transformer uses that information efficiently.

\section{Introduction and Background}
In-context learning offers a way to adapt predictions to a new system without updating model weights at test time. Here, each cube is a new system with fixed but hidden physical parameters. The central question is whether a transformer can infer enough from past interactions to predict the outcome of a new one.

\paragraph{Related work (expand to at least eight papers).}
\begin{enumerate}
    \item \textbf{Garg et al., ``What Can Transformers Learn In-Context?''} Study transformers learning function classes from input--output examples, providing a model-training and evaluation template. \url{https://arxiv.org/abs/2208.01066}
    \item \textbf{[Paper 2: Forgione et al.; confirm exact title/citation].} [At least one sentence summarizing the contribution.]
    \item \textbf{[Paper 3: CAPTURE; confirm exact citation].} [At least one sentence summarizing the contribution.]
    \item \textbf{[Paper 4].} [At least one sentence summarizing the contribution.]
    \item \textbf{[Paper 5].} [At least one sentence summarizing the contribution.]
    \item \textbf{[Paper 6].} [At least one sentence summarizing the contribution.]
    \item \textbf{[Paper 7].} [At least one sentence summarizing the contribution.]
    \item \textbf{[Paper 8].} [At least one sentence summarizing the contribution.]
\end{enumerate}

\section{Key Questions}
\begin{center}
\textbf{How much and what kind of context is informative enough for in-context learning?}
\end{center}
Specifically, how do the number and type of past pushes affect a transformer's ability to predict a new push on a cube with hidden mass $m$ and friction coefficient $\mu$, relative to an optimal estimator?
\begin{itemize}
    \item How does prediction error change with the number $c$ of context pushes? When does it approach the optimal estimator?
    \item For fixed $c$, how do varied forces, varied durations, and repeated pushes compare, particularly with noisy measurements?
    \item Does adding redundant pushes to an informative context hurt performance?
    \item Do these patterns generalize beyond the training ranges of $m$ and $\mu$?
\end{itemize}

\section{Hypotheses}
\begin{itemize}
    \item A small number of sufficiently informative pushes may let the transformer approach the optimal estimator. Since there are two unknown parameters ($m$ and $\mu$), we hypothesize that around two \emph{appropriately chosen} interactions could suffice in the ideal noise-free setting; parameter identifiability must be checked.
    \item Adding redundant pushes should not substantially increase prediction error.
    \item With measurement noise, repeated pushes should help average out noise even when they provide limited new parameter-identification information.
    \item We expect some of these trends to persist out of distribution, although absolute accuracy may decrease.
\end{itemize}

\section{Methods}
\subsection{Simulation and Dataset}
We use MuJoCo to simulate horizontal pushing of cubes with hidden mass $m$ and sliding friction $\mu$. Each cube provides a sequence of $C=11$ pushes, each defined by an applied force $F$ and duration $T$. We reset to a consistent initial state for each push and record the total displacement $d$, including movement during the applied force. We plan approximately 10,000 cubes, subject to compute availability.

\begin{itemize}
    \item Training: $m\in[0.2,2]$ kg (log-uniform), $\mu\in[0.2,0.8]$.
    \item OOD test: $m\in\text{[TODO]}$, $\mu\in\text{[TODO]}$.
    \item Pushes: $F\in\text{[TODO]}$ N, $T\in\text{[TODO]}$ s; choose $F>\mu m g$ so that each push initiates sliding.
    \item Measurement noise: $d_{\mathrm{obs}}=d+\varepsilon$, where $\varepsilon\sim\mathcal N(0,\sigma_\varepsilon^2)$ and $\sigma_\varepsilon\in\{0,\text{[TODO]},\text{[TODO]}\}$.
\end{itemize}

\subsection{Model and Training}
We train a small GPT-2-style causal transformer, following Garg et al., using alternating push and observed-displacement tokens. Each training sequence comes from one cube. For push $j$, let $x_j=(F_j,T_j)$ and let tildes denote normalized values. The model predicts the next displacement without seeing its label:
\begin{equation}
\hat d_j=f_\theta(\tilde x_1,\tilde d_{1,\mathrm{obs}},\ldots,\tilde x_{j-1},\tilde d_{j-1,\mathrm{obs}},\tilde x_j).
\end{equation}
We minimize mean squared error against the normalized \emph{noise-free} simulated displacement:
\begin{equation}
\mathcal L_b(\theta)=\frac{1}{NC}\sum_{n=1}^{N}\sum_{j=1}^{C}\left(\hat d_{n,j}^{(b)}-\tilde d_{n,j}^{(b)}\right)^2,
\end{equation}
where $N$ is the number of cube sequences per minibatch, $C=11$ is the number of pushes per sequence, and $b$ indexes minibatches. We use AdamW (learning rate, batch size, and training steps: \textbf{[TODO]}).

\subsection{Inference and Baselines}
At inference, the model receives $c$ observed push--displacement pairs from a single unseen cube, followed by a query push $(F,T)$; it outputs predicted total displacement $\hat d$. We evaluate separately for $c=0,\ldots,10$.
\begin{itemize}
    \item \textbf{No context:} prediction using only the query push ($c=0$).
    \item \textbf{Bayesian estimator:} grid over $(m,\mu)$ (resolution: \textbf{[TODO]}), with posterior weights proportional to the assumed parameter prior times the Gaussian likelihood of context measurements; predict the query using the posterior mean displacement. This is the squared-error-optimal estimator \emph{under the specified prior, simulator, and noise model}, up to grid approximation.
    \item \textbf{Oracle:} simulates the query using the cube's true $(m,\mu)$; for deterministic, noise-free targets its error is zero.
\end{itemize}

\subsection{Experiments and Evaluation}
\begin{enumerate}
    \item \textbf{Context amount:} error versus $c=0,\ldots,10$, compared with Bayesian and oracle baselines.
    \item \textbf{Diversity:} fix $c=\text{[TODO]}$; compare varied force, varied duration, and identical pushes at each noise level.
    \item \textbf{Redundancy:} append \textbf{[TODO]} repeated pushes to an informative context and measure the change in error.
    \item \textbf{OOD generalization:} repeat the above using cubes outside the training parameter ranges.
\end{enumerate}
We report $\mathrm{NMSE}(c)=\mathrm{MSE}(c)/\mathrm{MSE}(0)$, as well as absolute MSE, averaged over \textbf{[TODO]} random seeds with mean $\pm$ standard error. We also report training and validation curves. For fair comparisons, query distributions and evaluation targets will be held fixed.

\subsection{Implementation and Compute}
We use PyTorch, MuJoCo (\texttt{pip install mujoco}), and Garg et al.'s released transformer code (repository link: \textbf{[TODO]}). Our repository: \url{https://github.com/davidj24/cs182-final-project}.

\noindent\textbf{Hardware:} [TODO: HW9 P4 GPU, VRAM, CPU, RAM].\\
\textbf{Data generation:} [TODO: seconds/push] $\times$ [TODO: number of cubes] $\times$ 11 pushes.\\
\textbf{Training:} [TODO: steps, batch size, GPU-hours per run].\\
\textbf{Total budget:} [TODO: conditions $\times$ seeds $\times$ hours/run].

\clearpage
\appendix
\section{Generative-AI Use Disclosure}
\textbf{David Jo:} Used Gemini Notebook for an initial literature sweep; used Claude to scaffold the codebase and transcribe proposal notes into LaTeX.\\
\textbf{Soomi:} Used ChatGPT to brainstorm the project and evaluate feasibility. [Add any further assistance used.]\\
\textbf{Other team members:} [TODO].

\end{document}
